\chapter{Matrices}
\section{Matrices y sus Operaciones }
\begin{definition}[Matriz]
	Una matriz es un arreglo rectangular de números, estos números reciben el nombre de \textbf{entradas} o \textbf{elementos} de la matriz
\end{definition}
El tamaño de una matriz esta definida por el número de filas $m$ y el número de columnas $n$, además, es muy útil usar la notación de doble índice para referirse a las entradas de la matriz $a_{ij}$, donde el índice $i$ representa la i-ésima fila, y el ídice $j$ representa la j-ésima columna. Así una matriz de $m\times n$ se puede representar de la siguiente forma: 
\begin{equation}
	A = \left(  \begin{matrix}
		a_{11} & a_{12} &\cdots   &  a_{1n}\\
		a_{21}& a_{22} & \cdots  & a_{2n} \\
		\vdots &\vdots   &  \ddots & \vdots  \\
		a_{m1} & a_{m2} & \cdots  & a_{mn}
	\end{matrix}\right)
\end{equation}

\subsection{Suma de Matrices}
La suma de matrices se define por componentes, es decir, sea $A = (a_{ij})_{m\times n}$ y $B=(b_{ij})_{m\times n}$, entonces:
\begin{equation}
	A + B = (a_{ij}+ b_{ij})_{m \times n}
\end{equation}
es decir:
\begin{equation*}
		A +	B= \left(  \begin{matrix}
		a_{11} & a_{12} &\cdots   &  a_{1n}\\
		a_{21}& a_{22} & \cdots  & a_{2n} \\
		\vdots &\vdots   &  \ddots & \vdots  \\
		a_{m1} & a_{m2} & \cdots  & a_{mn}
	\end{matrix}\right) +  \left(  \begin{matrix}
	b_{11} & b_{12} &\cdots   &  b_{1n}\\
	b_{21}& b_{22} & \cdots  & b_{2n} \\
	\vdots &\vdots   &  \ddots & \vdots  \\
	b_{m1} & b_{m2} & \cdots  & b_{mn}
	\end{matrix}\right) = \left(  \begin{matrix}
	a_{11}+b_{11} & a_{12}+b_{12} &\cdots   &  a_{1n}+b_{1n}\\
	a_{21}+b_{21}& a_{22}+b_{22} & \cdots  & a_{2n}+b_{2n} \\
	\vdots &\vdots   &  \ddots & \vdots  \\
	a_{m1}+b_{m1} & a_{m2}+b_{m2} & \cdots  & a_{mn}+b_{mn}
	\end{matrix}\right)
\end{equation*}
\importante{NOTA: Para que la suma este definida, es necesario que tanto A como B sean del mismo tamaño}

\subsection{Multiplicación por un Escalar}
Sea A una matriz de $m \times n$ y $\alpha$ un escalar, la multiplicación por un escalar se define de la sigueinte forma:
\begin{equation}
	\alpha A = \alpha(a_{ij}) = (\alpha a_{ij}) 
\end{equation}
es decir:
\begin{equation*}
		\alpha A = \alpha \left(  \begin{matrix}
		a_{11} & a_{12} &\cdots   &  a_{1n}\\
		a_{21}& a_{22} & \cdots  & a_{2n} \\
		\vdots &\vdots   &  \ddots & \vdots  \\
		a_{m1} & a_{m2} & \cdots  & a_{mn}
	\end{matrix}\right) = \left(  \begin{matrix}
	\alpha a_{11} & \alpha a_{12} &\cdots   & \alpha a_{1n}\\
	\alpha a_{21}& \alpha a_{22} & \cdots  & \alpha a_{2n} \\
	\vdots &\vdots   &  \ddots & \vdots  \\
	\alpha a_{m1} &\alpha a_{m2} & \cdots  &\alpha a_{mn}
	\end{matrix}\right)
\end{equation*}
\importante{NOTA: A partir de estas definiciones, se puede trabajar con el negativo de una matriz:
$$
	A - B = A +(-B)
$$
 }
\subsection{Multiplicación de Matrices}
Si $A$ es una matriz de $m \times n$ y $B$ es una matriz $n \times r$, entonces, la multiplicación de matrices $C=AB$ es una matriz de $m \times r$, la entreda $c_{ij}$ se puede calcular con:
\begin{equation}
	c_{ij} = \sum^{n}_{k=1} a_{ik}b_{kj} = a_{i1}b_{1j}+a_{i2}b_{2j}+\dots+a_{in}b_{nj}
\end{equation}
visualmente se puede ver como la entrada $c_{ij}$ se produce de multiplicar la i-ésima fila con la j-ésima columna (como si fuera un producto punto)



\begin{equation*}
	C = AB = 
	\begin{bNiceMatrix}[
		margin,
		code-after = {
			% Resalta la fila 3 (i-ésima) de la primera matriz
			\tikz \node [fill=blue!20, fill opacity=0.5, inner sep=4pt, fit=(3-1)(3-4)] {} ;
		}
		]
		a_{11} & a_{12} & \dots  & a_{1n} \\
		\vdots & \vdots & \dots  & \vdots \\
		a_{i1} & a_{i2} & \dots  & a_{in} \\
		\vdots & \vdots & \ddots & \vdots \\
		a_{m1} & a_{m2} & \dots  & a_{mn}
	\end{bNiceMatrix}
	\begin{bNiceMatrix}[
		margin,
		code-after = {
			% Resalta la columna 3 (j-ésima) de la segunda matriz
			\tikz \node [fill=red!20, fill opacity=0.5, inner sep=0pt, fit=(1-3)(4-3)] {} ;
		}
		]
		b_{11} & \dots & b_{1j} & \dots & b_{1r} \\
		b_{21} & \dots & b_{2j} & \dots & b_{2r} \\
		\vdots & \vdots& \vdots & \dots & \vdots \\
		b_{n1} & \dots & b_{nj} & \dots & b_{nr}
	\end{bNiceMatrix}
\end{equation*}

\subsection{Potencia de Matrices:}
Sea A una matriz cuadrada. Se define $A^{k}$ como:
\begin{equation}
	A^{k}= A A \dots A 
\end{equation}

donde A se multiplica k-veces, con $k>0$.
\medspace
Es importante definir que cuando $k=0$ tenemos que:
\begin{equation}
	A^{0}= I_n
\end{equation}

\subsection{Transpuesta de una Matriz}
La transpuesta de una matriz $A$ de tamaño $m \times n$ se obtiene de intercambiar las filas por las columnas, se denota por $A^t$ y se define por:
\begin{equation}
	(A^t)_{ij}=A_{ji} \; \mathtt{  para \; toda }\; i\; \mathtt{ y \; } j
\end{equation}
